\section{Preliminaries}
\label{sec:bg}

We use the following notation: $\Sigma$ is a set of labels, $\Gamma$ is
a set of property names, and $\mathcal{D}$ is a (possibly infinite)
domain of property values.\looseness=-1

\stitle{Property Graphs}. %%
A property graph is a tuple $G=(V,E,\lambda,\rho)$, where $V$ is a
finite vertex set, $E \subseteq V \times V$ is a set of \emph{undirected} edges,
$\lambda: V \cup E \rightarrow \Sigma$ assigns labels to vertices and
edges, and $\rho: (V \cup E)\times \Gamma \rightarrow \mathcal{D}$ is a
partial property function. %
We use $x.\kw{prop}$ for the value of property $\kw{prop}$ on a vertex
or edge $x$.

\stitle{Subgraph Matching Queries}. A subgraph matching query is represented by a pattern
graph $Q_S=(V_Q,E_Q,\lambda_Q)$, where $V_Q$ and $E_Q$ are the query
vertex and edge sets, and $\lambda_Q$ is the query labeling function. A
\emph{subgraph match} of $Q_S$ in $G$ is an injective mapping
$M:V_Q \rightarrow V$ such that $\lambda_Q(u)=\lambda(M(u))$ for all
$u \in V_Q$, and for each $e=(u,v)\in E_Q$, there exists
$e'=(M(u),M(v))\in E$ with $\lambda_Q(e)=\lambda(e')$. We denote the set
of all such matches by $Q_S(G)$. %%
A \emph{path query} is a special subgraph query whose pattern graph is a
labeled path. We denote a path query by $P$ and a path query with $k$
edges is written as $P=u_0e_1u_1\cdots e_ku_k$, where $e_i$ connects
$u_{i-1}$ and $u_i$ for $i=1,\ldots,k$. We call $u_0$ and $u_k$ the
endpoints  of $P$.\looseness=-1


\eat{ %%
\stitle{Path Degrees}. %%
A \emph{path query} is a subgraph query whose pattern is a labeled path.
We denote it by $P=u_0e_1u_1\cdots e_ku_k$ and write $P(G)$ for its
match set in $G$. We refer to  $s=u_0$ as the source node of $P$ and
$t=u_k$ as the target node. 
For each vertex $v\in V$, define
$d_s^P(v)=\bigl|\{M\in P(G)\mid M(s)=v\}\bigr|$ and
$d_t^P(v)=\bigl|\{M\in P(G)\mid M(t)=v\}\bigr|$,
called the source and target \emph{path degrees} of $v$ \wrt $P$, respectively.
For a path degree $d>0$, we write $\operatorname{rank}(d)=k$ if $2^{k-1}<d\le 2^k$. %%
}%% 

\stitle{Graph Queries.}
A graph query augments a subgraph matching query with constraints on properties
 of query vertices and edges. Formally, a graph query is  a pair
 $Q=(Q_S,\psi)$, where $Q_S=(V_Q,E_Q,\lambda_Q)$ is a subgraph matching
 query and
 $\psi$ is a Boolean formula over literals on query vertices and edges.
 %%
A \emph{vertex literal} is an atomic condition on a query vertex $u\in
V_Q$ (\eg~$u.\kw{age}>30$), and an \emph{edge literal} is an atomic
condition on a query edge $e\in E_Q$ (\eg~$e.\kw{timestamp}<2025$). We
call both predicate literals. %%zh
A match of $Q$ in $G$ is any $M\in Q_S(G)$ with $M\models\psi$. For a
predicate literal $\ell$, $M\models\ell$ iff $\ell$ evaluates to true
under $M$; for vertex literal on $u$, evaluate on $M(u)$; for edge
literal on $e=(u,v)$, evaluate on $(M(u),M(v))$; $\neg,\land,\lor$ are
interpreted in standard semantics. %%
Hence, $Q(G)=\{M\in Q_S(G)\mid M\models\psi\}$.

\stitle{Cardinality Estimation.} Given a graph $G$ and a graph query
$Q$,  the \emph{cardinality estimation problem} aims to  compute an
estimate $\hat{C}$ of $|Q(G)|$ without explicitly computeing $Q(G)$. %


\begin{example}
\todo{Add a running example graph and query to illustrate the 
definitions and notations. @Xingzi}
\end{example}

\section{Path Degree Summary}

In this section, we present {path degree summary}, a core
summary data structure of our cardinality estimation framework. 

\vspace{.25em}%% 
We start with path degrees, which capture the exact distribution of
path counts for each vertex.  Unless stated otherwise, all summaries are
defined \wrt a fixed data graph $G=(V,E,\lambda,\rho)$. 

\stitle{Path Degrees.} %
Let $s$ and $t$ be the endpoints of a path query $P$. For each vertex
$v\in V$, define $d_s^P(v)=\bigl|\{M\in P(G)\mid M(s)=v\}\bigr|$ and
$d_t^P(v)=\bigl|\{M\in P(G)\mid M(t)=v\}\bigr|$, as the \emph{path
degrees} of $v$ \wrt $P$, respectively. %
The \emph{path degrees} of $P$ are the pair
$(D_s^P,D_t^P)$, where
\[D_s^P=\{(v,d_s^P(v))\mid d_s^P(v)>0\}, \quad D_t^P=\{(v,d_t^P(v))\mid
d_t^P(v)>0\}.\] %
For a degree value $d>0$, we write $\operatorname{rank}(d)=k$ if
$2^{k-1}<d\le 2^k$.

\stitle{Path Degree Summary.} %%
Directly using path degrees for cardinality estimation is
computationally and space expensive. %%
We therefore introduce PathDS, a compact summary of
path degrees that supports efficient estimation with low space
overhead.\looseness=-1

\begin{definitionx}\label{def:pathds} %%
Let $P$ be a path query with endpoints $s$ and $t$. The \emph{path
degree summary} (PathDS) of $P$ is a pair 
$(\mathcal{S}_s^P,\mathcal{S}_t^P)$, where 
\begin{mbi}%
\item $\mathcal{S}_s^P=\bigl\{(k,c)\mid
c=\lvert\{(v,d)\in D_s^P\mid
\operatorname{rank}(d){=}k\}\rvert,\ c>0\bigr\}$,
\item $\mathcal{S}_t^P=\bigl\{(k,c)\mid
c=\lvert\{(v,d)\in D_t^P\mid
\operatorname{rank}(d){=}k\}\rvert,\ c>0\bigr\}$. \qed\looseness=-1
\end{mbi}
\end{definitionx}

\vspace{.25em}%%
The PathDS summary $(\mathcal{S}_s^P,\mathcal{S}_t^P)$ can be
constructed from path degrees $(D_s^P, D_t^P)$ with a single linear
scan. %
Conceptually, $(\mathcal{S}_s^P,\mathcal{S}_t^P)$ is a rank-grouped
histogram of $(D_s^P, D_t^P)$ that groups path degrees by rank and
records the number of vertices in each rank. %
This design preserves the degree distribution at rank granularity while
substantially reducing space overhead.

\begin{example}
\todo{illustrate the definitions with a running example}
\end{example}

PathDS admits a $2$-approximation guarantee: replacing each degree in
rank $k$ by $2^k$ yields an estimate in $[|P(G)|,\,2|P(G)|]$.

\begin{propx} Let $P$ be a path query with two endpoints $s$ and $t$. Then\looseness=-1
\begin{mbi}
\item $|P(G)| \leq \sum \{2^k
\cdot c \mid {(k,c)\in \mathcal{S}_s^P} \}\leq 2|P(G)|$, and 
\item $|P(G)| \leq \sum \{2^k
\cdot c \mid {(k,c)\in \mathcal{S}_t^P} \}\leq 2|P(G)|$. \qed
\end{mbi}
\end{propx}

\stitle{Remark}. %\todo{Compare PathDS with PathCE, SafeBound}
Degree statistics have been widely used in cardinality estimation for
relational and graph queries,
\eg~\cite{deedsSafeBoundPracticalSystem2023,
pathce,caiPessimisticCardinalityEstimation2019,glogs}. The most related
methods are SafeBound~\cite{deedsSafeBoundPracticalSystem2023} and
PathCE~\cite{pathce}. We highlight their key differences from PathDS as
follows.

\sstab (a)  %% SafeBound vs. PathDS
SafeBound is designed for relational DBMS and uses ordered degree
sequences on join attributes,  which are essentially edge-level
statistics in graph terms. %%
PathDS differs from SafeBound in two aspects: (i) it focuses on
path-level degree statistics in contrast to edge-level statistics,  (ii)
it summarizes degree distributions into rank buckets, which are easier
to compute and maintain than the ordered degree sequences in
SafeBound (see also Exp.~\tbf).

\sstab (b) %% PathCE vs. PathDS
Both PathCE and PathDS use path-level statistics, but they summarize
them differently: (i) PathCE partitions vertices into summary groups
and stores maximal-degree statistics for each group pair; (ii) PathDS
groups path degrees by rank and stores per-rank counts, preserving
rank-level distributions with lower space overhead.\looseness=-1

\iffalse
\subsection{Graphs and Queries}

\noindent\textbf{Data Graph.} A data graph is defined as a triple $G = (V, E, L)$, where $V$ is a finite set
of vertices, $E \subseteq V \times V$ is a set of directed edges, and
$L$ is a labeling function that assigns each vertex $v \in V$ and edge
$e \in E$ a label together with a set of attributes
$(a_1, a_2, \ldots)$ from a domain $\mathcal{D}$.
This model subsumes commonly used property graph representations.

\noindent\textbf{Graph Query.} 
A graph query is represented as a pattern graph
$Q = (V_Q, E_Q, P)$, where $V_Q$ is a finite set of \emph{query vertices}
and $E_Q \subseteq V_Q \times V_Q$ is a set of \emph{query edges}.
Each query vertex represents a placeholder to be matched to a data
vertex in the input graph.
Each query edge $e = (u, v) \in E_Q$ specifies a structural constraint
between two query vertices, which may correspond to a single edge,
a fixed-length path, or a regular path expression (RPE).
The set $P = P_V \cup P_E$ consists of selection predicates on query
vertices and query edges, respectively.
For each $u \in V_Q$, predicates in $P_u \subseteq P_V$ restrict the
attributes of its matched data vertex, and for each $e \in E_Q$,
predicates in $P_e \subseteq P_E$ restrict the attributes of edges or
paths satisfying the structural constraint of $e$.

\noindent\textbf{Predicates and Selectivity.}
A predicate is a Boolean condition defined over vertex or edge attributes.
For a vertex predicate $P_u$, its selectivity is defined as the fraction
of vertices in $V$ that satisfy $P_u$.
Similarly, for an edge or path predicate $P_e$, its selectivity is the
fraction of candidate edges or paths satisfying the constraint.
Selectivity quantifies the filtering strength of predicates and is
fundamental to cardinality estimation.

\noindent\textbf{Graph Pattern Match.} Given a query $Q$ and a data graph $G$, a match of $Q$ in $G$ is a mapping $\rho: V_Q \rightarrow V$ such that: (i) for every variable $u \in V_Q$, $\rho(u)$ satisfies the predicates in $\mathcal{P}_u$; and (ii) for every query edge $e = (u, v) \in E_Q$, there exists a path in $G$ from $\rho(u)$ to $\rho(v)$ that conforms to the structural constraint of $e$ and satisfies the predicates in $\mathcal{P}_e$. We denote the set of all matches as $Q(G)$, assuming bag semantics for the results.


\subsection{PathSketch}

\noindent To support cardinality estimation of graph queries, $GCard$ maintains a data structure called \textit{PathSketch}, which summarizes statistics of path queries without predicate constraint. We define this data structure in two forms: the exact \textit{FlatDS} and its summary version, \textit{CompressedDS}.

\noindent\textbf{FlatDS.} Given a path pattern $P$ starting at variable $s$ and ending at $t$, let $V_s \subseteq V$ be the set of vertices matching $s$. The FlatDS of $P$ is defined as a sequence $f_P = (f_1, f_2, \dots, f_d)$, where $d = |V_s|$ and each element $f_i$ represents the number of distinct path instances of $P$ originating from the $i$-th vertex in $V_s$. Inspired by the degree sequence concept in SafeBound \cite{deedsSafeBoundPracticalSystem2023} and the short path statistics in PathCE\cite{full}, $f_P$ captures the exact frequency distribution of path endpoints. Notably, $f_P$ in our notation is an unordered list, where the elements are not required to be sorted by frequency.
\noindent\textbf{CompressedDS.}
To reduce storage overhead, the \textit{CompressedDS} represents
the frequency distribution of FlatDS as a set of compact tuples
$H = \{(\pi_1, \kappa_1), \dots, (\pi_d, \kappa_d)\}$.
For each frequency $f_i \in f_P$, we first extract an $n$-bit prefix
$\pi_i \in \{0,1\}^n$, corresponding to the most significant bits of
$f_i$ in its binary representation.
We then compute the bucket index
$\kappa_i = \lceil \log_b f_i \rceil$, which captures the order of
magnitude of $f_i$ under base $b$.
Together, $(\pi_i, \kappa_i)$ provides a compact representation
that preserves an upper bound on the original frequency while
significantly reducing space usage.

\eat{ %% EAT
\subsection{Cardinality Estimation}

\noindent\textbf{Problem Definition.} Following the standard definition in the field, the goal of cardinality estimation is to approximate the result size of a query without exhaustive execution \cite{harmouchCardinalityEstimationExperimental2017a, hanCardinalityEstimationDBMS2021}. Formally, given a data graph $G$ and a graph query $Q = (V_Q, E_Q, \mathcal{P})$, the problem is to compute an estimate $\hat{C}$ of the total number of matches $|Q(G)|$ . In the $GCard$ framework, the estimate $\hat{C}$ is derived by performing inference over PathSketch statistics while satisfying the structural constraints in $E_Q$ and the unified set of selection predicates $\mathcal{P}$. The objective is to produce $\hat{C}$ such that it remains as close as possible to the true cardinality $|Q(G)|$ while maintaining a negligible computational cost compared to the actual matching process.
} %% EAT
\fi
